\clearpage
\begin{examplebox}
\textbf{例1：原题 2021 全国甲卷（理）第20题，第一问怎样确定圆？}

抛物线顶点 $O$ 在原点、焦点在 $x$ 轴上，写作 $y^2=kx$（$k>0$）。直线 $x=1$ 与抛物线交于 $P=(1,\sqrt{k})$、$Q=(1,-\sqrt{k})$。由 $OP\perp OQ$，
\[
\overrightarrow{OP}\cdot\overrightarrow{OQ}
=1\times1+\sqrt{k}(-\sqrt{k})=1-k=0.
\]
所以 $k=1$，抛物线为 $\boxed{y^2=x}$。圆心 $M(2,0)$，与 $x=1$ 相切，所以半径 $r=|2-1|=1$，圆为
\[
\boxed{(x-2)^2+y^2=1}.
\]
\textbf{识别第二问：} $k=r=1$ 且 $h=2=1+1^2/1$，恰好满足 C096 主分支，故两条共顶点的弦切圆时，第三条弦也切圆。
\end{examplebox}

\clearpage
\begin{examplebox}
\textbf{例1（续）：不用同构、不用韦达，把第二问硬算完}

设 $A_i=(y_i^2,y_i)$。$A_iA_j$ 的方程是 $x-(y_i+y_j)y+y_iy_j=0$。
已知 $A_1A_2,A_1A_3$ 与圆相切，分别整理成
\[
\begin{cases}
(y_1^2-1)y_2^2+2y_1y_2+3-y_1^2=0,\\
(y_1^2-1)y_3^2+2y_1y_3+3-y_1^2=0.
\end{cases}
\]
\textbf{相减}并除以 $y_2-y_3\ne0$：
\[
(y_1^2-1)(y_2+y_3)=-2y_1.
\]
\textbf{相加}，用 $y_2^2+y_3^2=(y_2+y_3)^2-2y_2y_3$，利用刚求的和把一次项消去：
\[
(y_1^2-1)y_2y_3=3-y_1^2.
\]
从差式可知 $y_1^2\ne1$。因此
\[
2+y_2y_3=\frac{y_1^2+1}{y_1^2-1},\quad
1+(y_2+y_3)^2=\frac{(y_1^2+1)^2}{(y_1^2-1)^2}.
\]
圆心到 $A_2A_3$ 的距离平方是
\[
\frac{(2+y_2y_3)^2}{1+(y_2+y_3)^2}=1=r^2.
\]
所以 $A_2A_3$ 与圆相切。\textbf{整条路线都是普通硬算。}
\end{examplebox}

\clearpage
\begin{examplebox}
\textbf{例2：移动公共顶点，独立算出第三条直线}

仍取 $y^2=x$ 和 $(x-2)^2+y^2=1$，固定 $A_1=(1/4,1/2)$。从该点作两条切圆弦，其余交点的纵坐标 $b,c$ 满足
\[
-\frac34t^2+t+\frac{11}{4}=0
\quad\Longleftrightarrow\quad 3t^2-4t-11=0.
\]
所以 $b,c=(2\pm\sqrt{37})/3$。直接相加、相乘得
\[
b+c=\frac43,\qquad bc=-\frac{11}{3}.
\]
第三条弦：
\[
A_2A_3:\ x-\frac43y-\frac{11}{3}=0
\quad\Longleftrightarrow\quad \boxed{3x-4y-11=0}.
\]
圆心 $(2,0)$ 到它的距离为 $|6-11|/5=1$，确实切圆。\textbf{对照第 02 节右图：}它不是固定直线 $x=3$，但相切的性质没变。
\end{examplebox}

\clearpage
\begin{examplebox}
\textbf{例3：参数推广及反例，分清“特殊圆心”}

\textbf{推广：}若 $C:y^2=2x$、圆半径 $r=2$，主分支要求
\[
h=2+\frac{2^2}{2}=4.
\]
从 $O$ 作圆 $(x-4)^2+y^2=4$ 的两条切线，再与抛物线相交。设第二交点 $B=(b^2/2,b)$，直线 $OB$ 为 $2x-by=0$，相切条件给出 $|8|/\sqrt{4+b^2}=2$，即 $b^2=12$。因此另外两点为 $(6,\pm2\sqrt3)$，其连线 $x=6$ 到圆心 $(4,0)$ 的距离是 $2$，相切。

\medskip
\textbf{反例：} 若把原题圆改成 $(x-3)^2+y^2=1$，从 $O$ 向该圆作两条切线，第二交点为 $(8,\pm2\sqrt2)$，其连线 $x=8$ 到圆心的距离为 $5\ne1$。因此\textbf{普通轴上圆不能套 C096}。
\begin{center}
% 两个面板的坐标单位完全相同（各自 x=0...8.5, y=-3...3）。
\begin{tikzpicture}[x=.50cm,y=.50cm,>=Stealth,font=\scriptsize]
\begin{scope}
\draw[gray!50,->] (-.3,0)--(8.5,0) node[right] {$x$};
\draw[gray!50,->] (0,-3.15)--(0,3.15) node[above] {$y$};
\draw[blue!75!black,thick,domain=-2.85:2.85,samples=120,variable=\t] plot ({\t*\t},{\t});
\draw[teal!75!black,thick,fill=teal!8] (2,0) circle(1);
\draw[orange!85!black,very thick] (0,0)--(3,{sqrt(3)})--(3,{-sqrt(3)})--cycle;
\fill[teal!80!black](2,0)circle(1.8pt) node[below]{$M$};
\fill[orange!85!black](0,0)circle(1.5pt) node[above left]{$O$};
\node[font=\small] at (4,3.55) {$h=2$: 第三条切圆};
\node[orange!80!black,right] at (3.5,1.35) {$x=3$};
\end{scope}
\begin{scope}[xshift=5.7cm]
\draw[gray!50,->] (-.3,0)--(8.5,0) node[right] {$x$};
\draw[gray!50,->] (0,-3.15)--(0,3.15) node[above] {$y$};
\draw[blue!75!black,thick,domain=-2.85:2.85,samples=120,variable=\t] plot ({\t*\t},{\t});
\draw[teal!75!black,thick,fill=teal!8] (3,0) circle(1);
\draw[orange!85!black,very thick] (0,0)--(8,{sqrt(8)})--(8,{-sqrt(8)})--cycle;
\fill[teal!80!black](3,0)circle(1.8pt) node[below]{$M$};
\fill[orange!85!black](0,0)circle(1.5pt) node[above left]{$O$};
\node[font=\small] at (4,3.55) {$h=3$: 第三条不切圆};
\node[orange!80!black,right] at (8,1.8) {$x=8$};
\end{scope}
\end{tikzpicture}

\end{center}
\end{examplebox}
